% GENERATED FILE. Edit canonical lessons, then run npm run generate:books.
% canonical-source-hash: fnv1a64:88c12e7c336d8ec9
% canonical-lessons: ML-C223-cumakkuka, ML-C223-vaykkuka, ML-C223-dharikkuka, ML-C223-uruka, ML-C223-marruka

\chapter{To carry, To put, To wear, To take off, To change}
\label{ch:ml-to-carry-to-put-to-wear-to-take-off-to-change}
\hlchaptermodality{223}

\begin{chapteropening}
\textbf{By the end of this chapter:} \emph{I can say to carry, to put, to wear, to take off and to change.}

Complete the last lesson of chapter 223: I can say to carry, to put, to wear, to take off and to change.
\end{chapteropening}

\section[\texorpdfstring{\ml{ചുമക്കുക} (cumakkuka)}{ചുമക്കുക (cumakkuka)}]{\ml{ചുമക്കുക} (cumakkuka) — to carry}

\label{lesson:ML-C223-cumakkuka}



\begin{quote}
Before the new one: say the Malayalam for a barber, then the Malayalam for a tailor.
\end{quote}

\subsection*{You'll want to know: \ml{ചുമക്കുക}}
\textbf{\ml{ചുമക്കുക}} — \emph{cumakkuka} — \textquotedblleft{}to carry\textquotedblright{}.

\begin{cousinweb}
\textbf{\ml{ചുമക്കുക}} (\emph{cumakkuka}) is native Dravidian, one word with Tamil \emph{cuma} (DEDR 2677).
\end{cousinweb}

\begin{grammarlens}[title={What you've built}]
One more word for this chapter.
\end{grammarlens}

\subsection*{Your turn}
\begin{itemize}
\raggedright
  \item \emph{Say it:} \emph{cumakkuka}
  \item \emph{Say it:} \emph{cumakkuka}, once more
  \item \emph{Say it:} say \emph{cumakkuka}
  \item \emph{Recall:} say the Malayalam for a barber, then the Malayalam for a tailor, then say \emph{cumakkuka} again
\end{itemize}

\begin{tcolorbox}[breakable,colback=teal!4,colframe=teal!35!black,title={Before you move on}]
What does \ml{ചുമക്കുക} mean? (\textquotedblleft{}To carry\textquotedblright{}.) Say it once more.
\end{tcolorbox}
\section[\texorpdfstring{\ml{വയ്ക്കുക} (vaykkuka)}{വയ്ക്കുക (vaykkuka)}]{\ml{വയ്ക്കുക} (vaykkuka) — to put, to place}

\label{lesson:ML-C223-vaykkuka}



\begin{quote}
Before the new one: say the Malayalam for a tailor, then the Malayalam for to carry.
\end{quote}

\subsection*{You'll want to know: \ml{വയ്ക്കുക}}
\textbf{\ml{വയ്ക്കുക}} — \emph{vaykkuka} — \textquotedblleft{}to put, to place\textquotedblright{}.

\begin{cousinweb}
\textbf{\ml{വയ്ക്കുക}} (\emph{vaykkuka}) is native Dravidian, one word with Tamil \emph{vai}, Kannada \emph{bay} and Telugu \emph{vaicu} (DEDR 5549).
\end{cousinweb}

\begin{grammarlens}[title={What you've built}]
One more word for this chapter.
\end{grammarlens}

\subsection*{Your turn}
\begin{itemize}
\raggedright
  \item \emph{Say it:} \emph{vaykkuka}
  \item \emph{Say it:} \emph{vaykkuka}, once more
  \item \emph{Say it:} say \emph{vaykkuka}
  \item \emph{Recall:} say the Malayalam for a tailor, then the Malayalam for to carry, then say \emph{vaykkuka} again
\end{itemize}

\begin{tcolorbox}[breakable,colback=teal!4,colframe=teal!35!black,title={Before you move on}]
What does \ml{വയ്ക്കുക} mean? (\textquotedblleft{}To put, to place\textquotedblright{}.) Say it once more.
\end{tcolorbox}
\section[\texorpdfstring{\ml{ധരിക്കുക} (dharikkuka)}{ധരിക്കുക (dharikkuka)}]{\ml{ധരിക്കുക} (dharikkuka) — to wear}

\label{lesson:ML-C223-dharikkuka}



\begin{quote}
Before the new one: say the Malayalam for to carry, then the Malayalam for to put.
\end{quote}

\subsection*{You'll want to know: \ml{ധരിക്കുക}}
\textbf{\ml{ധരിക്കുക}} — \emph{dharikkuka} — \textquotedblleft{}to wear\textquotedblright{}.

\begin{cousinweb}
\textbf{\ml{ധരിക്കുക}} (\emph{dharikkuka}) is built on the Sanskrit root \emph{dhṛ}, \textquotedblleft{}to hold, to bear\textquotedblright{}, as in \emph{dhārayati}, \textquotedblleft{}holds, carries, wears\textquotedblright{} (CDIAL 6791), with the ending that makes a Malayalam verb.
\end{cousinweb}

\begin{grammarlens}[title={What you've built}]
One more word for this chapter.
\end{grammarlens}

\subsection*{Your turn}
\begin{itemize}
\raggedright
  \item \emph{Say it:} \emph{dharikkuka}
  \item \emph{Say it:} \emph{dharikkuka}, once more
  \item \emph{Say it:} say \emph{dharikkuka}
  \item \emph{Recall:} say the Malayalam for to carry, then the Malayalam for to put, then say \emph{dharikkuka} again
\end{itemize}

\begin{tcolorbox}[breakable,colback=teal!4,colframe=teal!35!black,title={Before you move on}]
What does \ml{ധരിക്കുക} mean? (\textquotedblleft{}To wear\textquotedblright{}.) Say it once more.
\end{tcolorbox}
\section[\texorpdfstring{\ml{ഊരുക} (ūruka)}{ഊരുക (ūruka)}]{\ml{ഊരുക} (ūruka) — to take off (clothes, shoes)}

\label{lesson:ML-C223-uruka}



\begin{quote}
Before the new one: say the Malayalam for to put, then the Malayalam for to wear.
\end{quote}

\subsection*{You'll want to know: \ml{ഊരുക}}
\textbf{\ml{ഊരുക}} — \emph{ūruka} — \textquotedblleft{}to take off (clothes, shoes)\textquotedblright{}.

\begin{cousinweb}
\textbf{\ml{ഊരുക}} (\emph{ūruka}) is native Dravidian, one word with Tamil \emph{ūr} (DEDR 750). Its first sense is \textquotedblleft{}to slip off\textquotedblright{}.
\end{cousinweb}

\begin{grammarlens}[title={What you've built}]
One more word for this chapter.
\end{grammarlens}

\subsection*{Your turn}
\begin{itemize}
\raggedright
  \item \emph{Say it:} \emph{ūruka}
  \item \emph{Say it:} \emph{ūruka}, once more
  \item \emph{Say it:} say \emph{ūruka}
  \item \emph{Recall:} say the Malayalam for to put, then the Malayalam for to wear, then say \emph{ūruka} again
\end{itemize}

\begin{tcolorbox}[breakable,colback=teal!4,colframe=teal!35!black,title={Before you move on}]
What does \ml{ഊരുക} mean? (\textquotedblleft{}To take off (clothes, shoes)\textquotedblright{}.) Say it once more.
\end{tcolorbox}
\section[\texorpdfstring{\ml{മാറ്റുക} (māṟṟuka)}{മാറ്റുക (māṟṟuka)}]{\ml{മാറ്റുക} (māṟṟuka) — to change}

\label{lesson:ML-C223-marruka}



\begin{quote}
Before the new one: say the Malayalam for to wear, then the Malayalam for to take off.
\end{quote}

\subsection*{You'll want to know: \ml{മാറ്റുക}}
\textbf{\ml{മാറ്റുക}} — \emph{māṟṟuka} — \textquotedblleft{}to change\textquotedblright{}.

\begin{cousinweb}
\textbf{\ml{മാറ്റുക}} (\emph{māṟṟuka}) is native Dravidian, the causative of the verb \textquotedblleft{}to be changed\textquotedblright{}; DEDR 4834 lists the root with Tamil \emph{māṟu}, Kannada \emph{māṟu} and Telugu \emph{māṟu}.
\end{cousinweb}

\begin{grammarlens}[title={What you've built}]
One more word for this chapter.
\end{grammarlens}

\subsection*{Your turn}
\begin{itemize}
\raggedright
  \item \emph{Say it:} \emph{māṟṟuka}
  \item \emph{Say it:} \emph{māṟṟuka}, once more
  \item \emph{Say it:} say \emph{māṟṟuka}
  \item \emph{Recall:} say the Malayalam for to wear, then the Malayalam for to take off, then say \emph{māṟṟuka} again
\end{itemize}

\begin{tcolorbox}[breakable,colback=teal!4,colframe=teal!35!black,title={Before you move on}]
What does \ml{മാറ്റുക} mean? (\textquotedblleft{}To change\textquotedblright{}.) Say it once more.
\end{tcolorbox}
